\section{Introduction}

Ferroelectric materials have extensive applications to advanced functional components, including capacitors,  non-volatile memories and ultrasonic transducers \cite{martin_thin-film_2017}. Currently, most of the widely used ferroelectric materials like \ce{Pb[Zr_{(1-x)} Ti_x]O3} (PZT) ceramics and \ce{Pb(Mg1/3 Nb2/3)O3-PbTiO3} (PMN-PT) crystals are intrinsically brittle, with the strength generally less than 100 MPa, and the fracture toughness merely around {1 $\mathrm{MPa}\cdot\mathrm{m}^{1/2}$} \cite{schneider_influence_2007}. However, as the core of actuators and ultrasonic generators,  ferroelectric materials are often subjected to mechanical vibrations and high-frequency electric fields. Such a severe environment can easily induce the initiation and propagation of cracks inside the ferroelectric materials. Accordingly, a comprehensive study on the fracture behaviors of ferroelectric materials is necessary for the reliable design of ferroelectric devices. 

The initiation and propagation of crack in ferroelectric material is accompanied by the microstructure evolution of polarization domain morphologies \cite{yingwei_influence_2020}. Such an electromechanical coupling property makes the fracture behaviors of ferroelectric materials rather complex. Thus, numerous experimental tests have been carried out for a better understanding of the fracture mechanism. For example, the fracture toughness of poled ferroelectric materials exhibits a significant anisotropic property in the Vickers indentation test \cite{yamamoto_internal_1983,okazaki_electro-mechanical_1992,lynch_crack_1995}. With the utilization of a high-speed camera, Wang et al. \cite{wang_influence_2020} captured the process of crack propagation in PZT-5H samples during the three-point bending fracture experiment. 

Strain gradient exist around the crack tip due to the formation of a singular stress field, leading to a significant flexoelectric effect \cite{zubko_flexoelectric_2013,yudin_fundamentals_2013}. Since ferroelectric materials with higher dielectric constants have large flexoelectric coefficients \cite{ma_flexoelectricity_2006,hong_first-principles_2013}, the crack tip is subjected to a giant strain gradient (flexoelectric) field, leading to a complex evolution of domain structure during the crack propagation process. For example, Wang et al. \cite{wang_direct_2020} observed a large polarization response near the crack tip in \ce{SrTiO3} epitaxial thin film induced by flexoelectric effect. Recently, Edwards et al. \cite{cordero-edwards_flexoelectric_2019} found that the flexoelectric effect around the crack tip can induce anisotropy in the fracture behavior of ferroelectric materials, and the crack propagation path shows asymmetrical characteristics.  Using an atomic force microscopy, Xu et al. \cite{xu_directly_2023} determined the electric field distribution around the crack tip, which provides a reliable approach to measuring the flexoelectric effect in ferroelectric materials. 

To gain a thorough understanding of the microscopic fracture mechanism in ferroelectric materials, the in situ observation and capturing of electromechanical response is indispensable. Currently, these requirements are still challenging for experimental investigation. Hence, numerical approach can further enrich the understanding of the fracture behavior in ferroelectric materials as complements to the experimental method. In recent years, the phase-field model has gained much attention in the study of microstructure evolution. The phase-field model utilizes a diffuse interface instead of a sharp one, which eliminates the requirements to track the moving of interface. Thus, the computation is greatly simplified \cite{chen_phase-field_2002}. A phase-field model has been widely used to investigate the polarization evolution in ferroelectrics \cite{wang_phase-field_2004,su_continuum_2007,hong_stability_2017,yuan_defect-mediated_2018}. Wang et al. \cite{wang_phase_2007,wang_phase_2008,wang_three-dimensional_2009,wang_effect_2010} utilized the phase-field to investigate the impact of domain switching on fracture toughness with different initial polarization directions, electric fields, and electrical boundary conditions. Using the phase-field model, Yu \cite{yu_i-integral_2016} et al. developed an I-integral method to analyze the crack-tip intensity factor, which is independent of domain switching and is suitable for large-scale domain switching. Zhao et al. \cite{zhao_effect_2018} incorporated the flexoelectric effect into the phase-field model for ferroelectric materials, which suggests that the flexoelectric effect can induce an asymmetric distribution of domain structures near the crack tip.

With the combination of Griffith fracture theory, the phase-field model can also be used to capture the fracture behaviors in ferroelectric materials. In the context of the linear piezoelectric theory, Tan et al. \cite{tan_phase_2023,tan_phase_2022} employed the phase-field model to investigate the fracture process of multiferroic materials  and the fatigue fracture behavior of ferroelectric materials. Taking the polarization evolution into account, Xu et al. \cite{xu_fracture_2010} established a phase-field continuum model for ferroelectric materials with a damage variable. Abdollahi et al. \cite{abdollahi_phase-field_2011,abdollahi_phase-field_2011-1,abdollahi_phase-field_2012,abdollahi_three-dimensional_2014} conducted a series of phase-field models to study the coupling between crack propagation and domain evolution in ferroelectric single crystals and polycrystals. Zhang et al. \cite{yong_zhang_phase_2022,zhang_jumping_2022} introducing the dielectric breakdown as an additional order parameter in the phase-field model. The evolution of polarization, dielectric breakdown, and crack propagation can consequently be calculated simultaneously. Nevertheless, the flexoelectric effect has yet to be incorporated into the phase-field model for the fracture of ferroelectric materials to the best of the author's knowledge.

In this work, a phase-field model for the brittle fracture in ferroelectric materials with the consideration of the flexoelectric effect is developed. To satisfy the high-order continuity requirements of spatial discretization, the isogeometric analysis method is employed. Additionally, polynomial splines over hierarchical T-meshes (PHT-splines) are used to enable local mesh refinement near the crack, thereby enhancing computational efficiency. With this model, the crack evolution with a different initial polarization direction is simulated. The effect of polarization domain switching and flexoelectricity throughout the crack extension process is investigated. 

The structure of this paper is organized as follows. Section \ref{Section2} introduces the framework of the phase-field model and the corresponding governing equations. The numerical implementation is detailed in Section \ref{Section3}. Section \ref{Section4} presents the simulation results and discussions. The concluding remarks are draw in Section\ref{Section5}.